\section{An Evaluated Coupling Candidate and Its Scale Limits}
\label{v480:coupling-scale-candidate}

The original field equation names effective coefficients and triadic
operators but does not evaluate their joint action. This section
therefore does two separate things: it evaluates a declared cyclic
candidate using the inherited phase input, and tests whether its
remaining scales could produce the required mass ratios. It does
not identify that candidate with the native clock cross block.

\subsection{Separating Shape, Normalization and Physical Scale}
\label{v480:coupling-normalization}

For a nonzero supplied cross block $V$ on a finite positive-metric
realization, the convention $\operatorname{Tr}(B^*B)=1$ fixes
$\mu=\|V\|_{\rm F}$ and $B=V/\mu$. It removes a redundancy,
not an undetermined physical coupling.

\begin{proposition}[Coupling and gap scales remain independent input]
\label{v480:coupling-scale-freedom}
In $V=\mu B$, replacing $(\mu,B)$ by $(t\mu,B/t)$, $t>0$,
leaves $V$ unchanged. Normalizing $B$ removes this redundancy.
For normalized $B$ in the flat-gap model of
Proposition~\ref{v480:native-flat-gap-mass}, put
$\eta=\mu/\Delta$. Then
\[
 m_i=\Delta\,\frac{\sqrt{1+4\eta^2q_i}-1}{2}.
\]
The shape $q_i$ and $\eta$ determine ratios, but do not fix
$\Delta$. Simultaneously scaling $(\mu,\Delta)$ by $s>0$
scales all masses by $s$, leaving ratios and complementary weights
unchanged.
\end{proposition}
\begin{proof}
Substitution proves the redundancy and dimensionless mass formula.
The exact complementary weight is $m_i/(\Delta+2m_i)$, which
is invariant under the simultaneous scaling.
\end{proof}

Even a supplied $Q$ leaves $B$ undetermined up to output isometries
in its polar decomposition; these do not change $B^*B$.
The source symbols $g_Y$, $\chi$ and the paired triadic operators
specify products, not a numerical factorization into a scale and a
normalized matrix. A physical normalization needs canonical kinetic
terms, common units and actual matrix elements.

\subsection{A Minimal Two-Channel Guess Is Insufficient}

\begin{proposition}[Two cyclic harmonics cannot give the Koide ratio]
\label{v480:coupling-two-harmonic-bound}
For the generation cycle $U$, let $B=b_0I+b_1U$, with arbitrary
complex $b_0,b_1$ not both zero. If the candidate masses are
proportional to the eigenvalues of $B^*B$, their Koide diagnostic
satisfies $K\leq1/2$, not $2/3$.
\end{proposition}
\begin{proof}
In the Fourier basis let $\omega=e^{2\pi i/3}$ and
$\lambda_k=b_0+b_1\omega^k$. The numbers
$z_k=\omega^k\lambda_k$ have zero sum, since both
$\sum_k\omega^k$ and $\sum_k\omega^{2k}$ vanish.
Their magnitudes are the principal root masses up to a common
factor. The closed-triole norm bound of
Proposition~\ref{v480:hc-mass-triole-bound} therefore gives
$\sum_k|z_k|^2/(\sum_k|z_k|)^2\leq1/2$.
\end{proof}

Thus a two-channel relative phase is not by itself enough in this
minimal cyclic realization. More general operator-valued channels
are not excluded by the proposition.

\subsection{Numerical Three-Channel Candidate From the Written Phase}
\label{v480:coupling-evaluated-matrix}

Take the \emph{inherited hypotheses}, not newly proved laws,
$r/a=1/\sqrt2$ and
$\theta=\pi-\sqrt2\arcsin(\bsym)$, with the paper's rounded
input $\bsym=0.5493$. Fix only the normalization
$\operatorname{Tr}(B^*B)=1$. On the positive-root phase branch,
\[
 \theta=2.319190257857\ldots,\qquad
 B=\frac{I_3}{\sqrt6}
   +\frac{e^{i\theta}U+e^{-i\theta}U^*}{\sqrt{12}}.
\]
In the convention $U\varepsilon_i=\varepsilon_{i+1\bmod3}$,
\[
 B\simeq
 \begin{pmatrix}
  0.408248&-0.196433-0.211536i&-0.196433+0.211536i\\
  -0.196433+0.211536i&0.408248&-0.196433-0.211536i\\
  -0.196433-0.211536i&-0.196433+0.211536i&0.408248
 \end{pmatrix}.
\]
Its normalized Gram matrix has diagonal $1/3$ and lower cyclic
coefficient $-0.166548+0.255824i$. Its ascending eigenvalues are
\[
 (q_e,q_\mu,q_\tau)
 \simeq(0.0002366350,\ 0.0567819436,\ 0.9429814214).
\]
The subscripts denote an attempted ascending lepton assignment,
not an independently derived generation dictionary. The leading
ratio prediction \emph{conditional on this shape} is
\[
 \frac{q_e}{q_\tau}=0.0002509434\ldots,\qquad
 \frac{q_\mu}{q_\tau}=0.0602153365\ldots.
\]
The inherited mass benchmark in Theorem~\ref{thm:koide-triolic},
used here only for comparison, instead gives approximately
$0.0002875860$ and $0.0594633229$. The relative ratio differences
are $-12.74\%$ and $+1.265\%$. No measured lepton mass enters
the construction of $B$. Nevertheless the radius, phase law and
rounded frame input are assumptions, so this is not a parameter-free
native derivation. The Koide diagnostic $2/3$ is inherited from the
imposed radius.
At $\bsym=0$ the same fixed-radius formula is not a positive root
matrix and does not give three degenerate masses. This local candidate
is not a valid interpolation to the proposed photon limit.

An explicit effective six-dimensional realization is
\[
 P=\begin{pmatrix}I_3&0\\0&0\end{pmatrix},\quad
 S=\begin{pmatrix}I_3\\0\end{pmatrix}.
\]
\[
 \widetilde R_{\rm cand}=\begin{pmatrix}0&B^*\\B&0\end{pmatrix},
 \quad
 \Omega_{\rm cand}=\begin{pmatrix}0&\mu B^*\\\mu B&\Delta I_3\end{pmatrix}.
\]
It realizes the candidate Gram matrix and field cross block, with
$V=\mu B$. It has not been shown to be the paper's
$R_{\mathrm{GUE}}$, a native Millennium realization or a
generation-compatible embedding of its Dixon carrier. Its imposed
Hermitian off-sector clock is not that missing identification.

\subsection{A Scale Choice Cannot Repair the Fixed Shape}

\begin{proposition}[The positive-gap correction increases smaller mass ratios]
\label{v480:coupling-gap-ratio-obstruction}
In the flat positive-gap model, for $0<q_i<q_j$ and $\mu>0$,
\[
 \frac{m_i}{m_j}
 =\frac{q_i}{q_j}
  \frac{\sqrt{\Delta^2+4\mu^2q_j}+\Delta}
       {\sqrt{\Delta^2+4\mu^2q_i}+\Delta}
 >\frac{q_i}{q_j}.
\]
Consequently no positive $\mu,\Delta$ can reproduce the quoted
muon/tau benchmark with the fixed numerical shape above.
\end{proposition}
\begin{proof}
Rationalize each exact mass from
Proposition~\ref{v480:native-flat-gap-mass}. The numerator of the
second factor is strictly larger than its denominator. Already
$q_\mu/q_\tau$ exceeds the quoted benchmark; the correction can
only increase that ratio. Common scaling cannot alter it.
\end{proof}

This excludes this fixed shape with zero direct block and flat
positive gap, not all coupling matrices, non-flat gaps or direct
internal mass terms. A change of the relative phase, radius or
field blocks needs an independent structural law; fitting it to
the benchmark is not such a derivation.

\subsection{Reference Scales and Explicit Calibration Boundaries}
\label{v480:coupling-scale-reference}

A reference mass $m_{\rm ref}$ permits the parametrization
$\mu=u\,m_{\rm ref}$, $\Delta=d\,m_{\rm ref}$ with $u,d>0$.
It does not fix the dimensionless coefficients $u,d$.
For comparison only, the ordinary single-frame Planck reference is
$m_{\rm P}=\sqrt{\hbar c/G}$, or $\hbar/(c\ell_{\rm P})$.
The external CODATA central reference is
$m_{\rm P}c^2=1.220890\times10^{19}\,\mathrm{GeV}$
\footnote{NIST, 2022 CODATA, Planck mass energy equivalent:
\url{https://physics.nist.gov/cgi-bin/cuu/Value?plkmc2gev}.}.
This external normalization is not predicted by the present construction.

As an independently defined electroweak reference, use
$v_F=(\sqrt2G_F)^{-1/2}=246.21965\ldots\,\mathrm{GeV}$ and
$\mu_{\rm EW}=v_F/\sqrt2=174.10358\ldots\,\mathrm{GeV}$,
with the CODATA central value
$G_F=1.1663787\times10^{-5}\,\mathrm{GeV}^{-2}$
\footnote{NIST, 2022 CODATA, Fermi coupling constant:
\url{https://physics.nist.gov/cgi-bin/cuu/Value?gf}.}.
Using this reference for $\mu$ would be an additional choice,
not a recovered Yokamono/Higgs coupling.

With $\Delta=m_{\rm P}$ and $\mu=\mu_{\rm EW}$, the leading
coefficient is $\kappa=2.48278\times10^{-6}\,\mathrm{eV}$ and
the exact positive masses are approximately
\[
 (5.8751\times10^{-10},\ 1.4098\times10^{-7},\
  2.3412\times10^{-6})\,\mathrm{eV}.
\]
Thus this particular Planck-gap/electroweak-mixing choice does not
give the charged-lepton scales quoted in the source. Choosing
both scales at the Planck reference instead gives masses of
Planck order for the larger candidate eigenvalues; neither
choice is a model-derived scale identification.

For clarity, a single mass can always be used as a calibration:
\[
 \mu^2=\frac{m_\tau(m_\tau+\Delta)}{q_\tau},\qquad
 \Delta=\frac{\mu^2q_\tau}{m_\tau}-m_\tau.
\]
With the inherited $m_\tau=1.77686\,\mathrm{GeV}$, fixing
$\Delta=m_{\rm P}$ gives $\mu\simeq4.7964\times10^9\,\mathrm{GeV}$;
fixing $\mu=\mu_{\rm EW}$ gives
$\Delta\simeq1.6085\times10^4\,\mathrm{GeV}$.
These are explicitly one-mass calibrations, not predictions, and
neither removes the muon/tau ratio obstruction.

The paper's proposed frame-insertion assignments must also be kept
consistent if used. Formally applying
$\hbar_f=\gamma_L^{-1}\hbar_m$, $G_f=\gamma_L^{-2}G_m$ and
$c_f=c_m$ gives $m_{{\rm P},f}=\gamma_L^{1/2}m_{{\rm P},m}$.
Using $\hbar_m/(c_m\ell_{{\rm P},f})$ instead would mix frame
assignments and introduce an extra factor $\gamma_L$.
This is conditional bookkeeping, not validation of the frame laws.

\subsection{Diagnosing the Matrix Discrepancy}
\label{v480:matrix-phase-diagnosis}

The matrix introduces no extra distortion of the inherited Brannen
spectrum: its Fourier eigenvalues are exactly those root coordinates,
up to the common normalization. Direct diagonalization and the
analytic formula agree. Six-decimal rounding of its displayed entries
changes the electron/tau ratio by about $0.01225\%$, not $12.74\%$.
Cyclic relabelling and complex conjugation only permute its spectrum.

For a diagnostic inversion of the \emph{historical comparison only},
put $x_k=\sqrt{m_k/\sum_jm_j}$ in the ascending assignment above.
Then
\[
 a_{\rm cmp}=\tfrac13\sum_kx_k,\qquad
 z_{\rm cmp}=\tfrac13\sum_kx_k e^{-2\pi i k/3},\qquad
 \rho_{\rm cmp}=|z_{\rm cmp}|/a_{\rm cmp}.
\]
The discrete inverse Fourier identity reconstructs these three positive
roots. It gives
\[
 \theta_{\rm cmp}=2.316624255394\ldots,\qquad
 \rho_{\rm cmp}=0.707100578472\ldots.
\]
Relative to the imposed $\rho=1/\sqrt2$, the radius difference is
about $-8.77$ parts per million, whereas the phase difference is
$-0.00256600246\ldots$ rad, or $-0.147021\ldots$ degrees.
Changing only the phase to this comparison-derived value gives
ratios approximately $(0.0002874631,0.0594619634,1)$; changing
only the radius leaves $(0.0002510584,0.0602167023,1)$.
Thus the phase accounts for almost all of this fixed-ansatz discrepancy.
These inversions are post-hoc diagnostics, not a new derived phase law.
This $\theta$ is the relative coefficient phase of $U$ and $U^*$,
not the common pacer phase, which still cancels from the Gram matrix.

\begin{proposition}[Relative-phase sensitivity near a small root]
\label{v480:matrix-phase-sensitivity}
Let $h_k=1+2\rho\cos\phi_k>0$, with
$\phi_k=\theta+2\pi k/3$ and
$q_k=h_k^2/(3+6\rho^2)$. For a fixed reference mode $j$,
\[
 \partial_\theta\log\frac{q_i}{q_j}
 =-4\rho\left(\frac{\sin\phi_i}{h_i}
                    -\frac{\sin\phi_j}{h_j}\right),
 \qquad
 \partial_\rho\log\frac{q_i}{q_j}
 =4\left(\frac{\cos\phi_i}{h_i}
                    -\frac{\cos\phi_j}{h_j}\right).
\]
For the inherited phase law,
$\partial_\beta\theta=-\sqrt2/\sqrt{1-\beta^2}$.
\end{proposition}
\begin{proof}
The common normalization cancels from $q_i/q_j$.
Differentiate $2\log h_i-2\log h_j$ and apply the chain rule.
\end{proof}

At the candidate, the electron/tau and muon/tau logarithmic derivatives
with respect to $\beta$ are approximately $92.6426$ and $-8.28621$.
The small electron root is close to the phase-chamber boundary, which
explains the large relative amplification. This is not evidence of an
unrepresented physical frame interaction by itself.

Within the stated phase law, the comparison phase would correspond to
$\beta_{\rm cmp}=0.5508152849\ldots$, a shift of $0.00151528$.
That is about thirty half-steps of four-decimal rounding. If decimal
rounding alone is tested on $[0.54925,0.54935]$, the middle/largest
ratio ranges from about $0.06024029$ down to $0.06019039$ and remains
above the comparison. Its monotonic decrease here follows from the
signs in the displayed derivative. This interval is a rounding test,
not an assigned experimental uncertainty.
Using the exact written $\gamma_{\rm string}=\log2/(\pi\sqrt3)$
with the still approximate $\gamma_{\rm BI}=0.2375$ gives
$\beta=0.5493482484\ldots$ and ratio $0.06019127\ldots$;
it does not remove the discrepancy. Inverting to a new
$\gamma_{\rm BI}\simeq0.2383454$ would be another calibration,
not a justified revision of that input.

\subsection{What the Planck Test Does and Does Not Exclude}
\label{v480:planck-assumption-diagnosis}

An independent SI reconstruction gives
$m_{\rm P}=2.17643434\ldots\times10^{-8}\,\mathrm{kg}$ and
$\ell_{\rm P}=1.61625502\ldots\times10^{-35}\,\mathrm{m}$.
Converting $m_{\rm P}c^2$ from joules to GeV reproduces
$1.22089013\ldots\times10^{19}\,\mathrm{GeV}$, consistent with
the rounded CODATA reference above
\footnote{The SI check uses the exact defining values of $h,c,e$
and the NIST CODATA central $G=6.67430\times10^{-11}$ in SI units:
\url{https://physics.nist.gov/cgi-bin/cuu/Value?bg}.}.
There is no missing GeV/eV factor or missing power of $c$ in the
micro-eV estimate.

The suppression comes from the three added assumptions
$A=0$, $\Delta=m_{\rm P}$ and $\mu=\mu_{\rm EW}$:
the leading energy is $\mu^2q_i/\Delta$. The largest value is
about $7.59\times10^{14}$ times below the quoted tau mass.
Substituting the reduced reference $m_{\rm P}/\sqrt{8\pi}$ as
the physical gap changes the result by only $\sqrt{8\pi}$,
to about $1.174\times10^{-5}\,\mathrm{eV}$ for the largest mass.
That substitution is another gap choice, not a unit correction.
The paper's conditional consistent Planck frame factor is only
$\gamma_L^{1/2}\simeq1.3654$; even the inconsistent mixed-$\hbar$
factor is only $\gamma_L^{3/2}\simeq2.5453$.
Neither accounts for the required hierarchy.

A minimum reference length or UV cutoff does not, without a spectral
equation, fix the complementary mass gap in this block model.
Keeping $\mu=\mu_{\rm EW}$ and calibrating the tau mass gives
$\Delta\simeq16.085\,\mathrm{TeV}$, corresponding to the inverse
energy length $\hbar c/\Delta\simeq1.227\times10^{-20}\,\mathrm m$.
This is about $7.59\times10^{14}\ell_{\rm P}$, an effective gap
length, not a new fundamental Planck length. Keeping the gap at the
Planck reference instead requires an intermediate mixing reference
$\mu\simeq4.7964\times10^9\,\mathrm{GeV}$, about
$2.755\times10^7\mu_{\rm EW}$, in that one-mass calibration.
Neither possibility is fixed by the source or repairs the flat-gap
ratio obstruction.

\begin{proposition}[Non-flat gaps can absorb a target spectrum]
\label{v480:planck-nonflat-fit-freedom}
Suppose the positive candidate is
$B=F\operatorname{diag}(\sqrt{q_i})F^*$, $q_i>0$.
For any target magnitudes $t_i>0$, choose $\mu$ such that
$\mu^2q_i>t_i^2$ and set
\[
 \Delta_i=\frac{\mu^2q_i}{t_i}-t_i>0,\qquad
 H_\perp=F\operatorname{diag}(\Delta_i)F^*.
\]
Then $\left(\begin{smallmatrix}0&\mu B^*\\\mu B&H_\perp
\end{smallmatrix}\right)$ has light signed eigenvalues $-t_i$.
\end{proposition}
\begin{proof}
Conjugating by $\operatorname{diag}(F,F)$ splits the system into
three $2\times2$ blocks. Their light magnitudes obey
$t_i(t_i+\Delta_i)=\mu^2q_i$, exactly the defining relation.
\end{proof}

For the historical comparison, fixing $\Delta_\tau=m_{\rm P}$
and the tau-calibrated intermediate $\mu$ gives
\[
 (\Delta_e,\Delta_\mu,\Delta_\tau)/m_{\rm P}
 \simeq(0.87258583,1.01264668,1).
\]
All gaps remain of Planck order and all three historical masses are
reproduced, by construction. This is an exact fitted counterexample
to a blanket exclusion of Planck-order gaps, not a prediction of these
anisotropies. It uses the target masses to define $H_\perp$.
It cannot determine whether a physical discrepancy belongs to the
root coupling, the gap or another term. With electroweak mixing
and every gap still of Planck order, this modest gap anisotropy
does not fix the absolute suppression.

\begin{proposition}[A retained direct block makes the Planck sector a small correction]
\label{v480:planck-direct-block-bound}
For the Hermitian block model of Proposition~\ref{v480:native-schur-mass},
let $H_\perp\geq\Delta I$ and $|\lambda|\leq E<\Delta$.
Every such eigenvalue satisfies
\[
 \operatorname{dist}(\lambda,\operatorname{spec}A)
 \leq\frac{\|V\|^2}{\Delta-E}.
\]
\end{proposition}
\begin{proof}
The closed coordinate $x$ cannot vanish, since otherwise the
eigenvalue would lie in $\operatorname{spec}H_\perp$.
The Schur equation gives
$(A-\lambda I)x=V^*(H_\perp-\lambda I)^{-1}Vx$.
The inverse norm is at most $(\Delta-E)^{-1}$. For self-adjoint
$A$, the spectral theorem bounds
$\|(A-\lambda I)x\|$ below by
$\operatorname{dist}(\lambda,\operatorname{spec}A)\|x\|$.
\end{proof}

With the candidate $\|V\|^2=\mu_{\rm EW}^2q_{\max}$,
$\Delta=m_{\rm P}$ and $E=2\,\mathrm{GeV}$, the bound is about
$2.3412\times10^{-6}\,\mathrm{eV}$. A direct block at charged-lepton
scales could therefore coexist with Planck-gap corrections that are
tiny. The original $\Omega$ already contains direct/projected/triadic
terms; setting their compressed sum $A$ to zero was an additional
choice of the earlier induced-mass test. This alternative does not
derive $A$ or its hierarchy, but it shows why that test does not
exclude the paper's general use of the Planck reference.

\paragraph{Determination reached.}
There is now an explicit numerical normalized coupling candidate
and a testable six-dimensional block realization. Its relative
shape is conditional on the inherited radius and phase hypotheses.
No numerical native values of $g_Y$, $\chi$, $\mu$ or $\Delta$
have been established. The missing physical input is a represented
cross-block interaction and gap equation that fixes them independently
of the target lepton masses, and resolves the strict-versus-dressed
closure boundary of Section~\ref{v480:native-operator-link}.
The discrepancy diagnosis separates the inherited relative-phase
ansatz from matrix arithmetic, and the extra flat-Planck-gap induced
mass hypothesis from unit conversion or a fundamental length claim.
It does not uniquely identify a physical cause from the mass data.
